\documentclass[11pt,a4paper]{article}
\usepackage[utf8]{inputenc}
\usepackage[spanish]{babel}
\usepackage{amsmath, amssymb}
\usepackage{booktabs}
\usepackage{geometry}
\geometry{margin=1in}
\title{Reporte Forense Ambiental y Atribución Atmosférica (Últimas 96 Horas)}
\author{Departamento de Análisis Forense Ambiental}
\date{\today}

\begin{document}
\maketitle

\section{Dictamen Termodinámico (Validación iPINN)}
El análisis de consistencia física (PVI: 0.025999) confirma una atmósfera estable con inversión térmica persistente. La divergencia del campo de velocidades es despreciable, validando la incompresibilidad del flujo.

\begin{table}[h]
\centering
\begin{tabular}{lc}
\toprule
\textbf{Métrica} & \textbf{Valor} \\
\midrule
Physics Violation Index (PVI) & 0.025999 \\
Divergencia Máxima Absoluta & 0.364332 \\
Estado de Estabilidad & Estable (Inversión Térmica) \\
\bottomrule
\end{tabular}
\caption{Indicadores de consistencia física.}
\end{table}

\section{Atribución Geoespacial de Emisiones (PM2.5)}
Se identificaron dos cúmulos críticos mediante clustering GMM:

\begin{itemize}
    \item \textbf{Cúmulo 0 (Ladera Occidental):} Tasa de emisión $S = 0.000307 \mu g/(m^3 s)$. Foco residencial y de alta pendiente.
    \item \textbf{Cúmulo 1 (Eje Central):} Tasa de emisión $S = 0.000234 \mu g/(m^3 s)$. Corredor del Río Medellín.
\end{itemize}

\section{Plan de Acción Adaptativo (PAA-96H)}
Ante la alerta Naranja-Roja, se establecen las siguientes medidas:

\subsection{Movilidad}
\begin{itemize}
    \item Extensión de Pico y Placa a 4 dígitos en el corredor del Río.
    \item Restricción total de carga pesada (pre-2015) de 05:00 a 20:00.
\end{itemize}

\subsection{Regulación Industrial}
\begin{itemize}
    \item Suspensión de procesos de combustión en zona de influencia del Cúmulo 0.
    \item Auditoría obligatoria de filtros HEPA en fuentes fijas de la ladera.
\end{itemize}

\end{document}